% !TeX root = ../main.tex

% ==================================================
% D.6. 불리언 마스크와 조건부 선택
% ==================================================

\subsection{불리언 마스크와 조건부 선택}

\subsubsection{이전 학습과의 연결}

\zsh{D.4}에서는 Pandas의 데이터를 선택하는 두 가지 좌표계인

\zsh{레이블(label)}과 \zsh{위치(position)}를 구분했다.

\zsh{D.5}에서는 한 단계 더 나아가, 서로 다른 Series와 DataFrame을 연산할 때 Index 레이블을 기준으로 값의 대응 관계가 결정된다는 사실을 배웠다.

이번 학습에서는 이 두 개념을 결합할 것이다.

다음 코드를 살펴보자.

\begin{dhkcode}
import pandas as pd

df = pd.DataFrame(
    {"score": [90, 70, 85]},
    index=["A", "B", "C"],
)

condition = df["score"] >= 80

result = df.loc[condition]
\end{dhkcode}

이 코드는 점수가 80점 이상인 행을 선택한다.

그런데 여기서 중요한 질문은 단순히 어떤 행이 반환되는지가 아니다.

\zsh{condition}은 어떤 객체이며, 그 객체가 \zsh{df.loc[...]}에 전달되었을 때 Pandas는 어떤 기준으로 행을 선택할까?

실행 과정을 나누어 보면 다음과 같다.

\begin{dhksmry}
    1. 비교식 평가
       df["score"] >= 80
       불리언 Series 생성 · Index 유지

                \da

    2. 인덱서에 조건 전달
       df.loc[condition]
       조건 Series의 레이블을 DataFrame의 Index에 대응

                \da

    3. 조건부 행 선택
       레이블 A, C에 해당하는 행 반환
       새로운 DataFrame 결과 · 원본은 변경되지 않음
\end{dhksmry}

여기서 가장 중요한 구분은 조건을 계산하는 과정과 조건을 적용하는 과정은 서로 다른 연산이라는 사실이다.

\subsubsection{이번 학습의 핵심 질문}

이번 학습에서는 다음 내용을 분석할 것이다.

\begin{dhksmry}
1. 불리언 마스크(Boolean mask)는 어떤 객체이며, 어떻게 생성되는가?

2. &, |, ~는 Python의 and, or, not과 어떻게 다른가?

3. 마스크에 <NA>가 존재하면 조건 계산과 행 선택에서 각각 어떻게 
   처리되는가?

4. isin()은 일반적인 비교 연산과 어떤 차이가 있는가?

5. where()와 mask()는 행 필터링과 어떻게 다른가?

6. query()는 어떤 방식으로 조건식을 평가하며, Python 코드 및 SQL과 
   무엇이 다른가?

7. 마스크의 Index가 선택 대상의 Index와 다르면 어떤 일이 발생하는가?

8. ML/Transformer 분석에서 특정 클래스의 오류 샘플과 길이 초과 문서를 
   어떻게 안전하게 추출하는가?
\end{dhksmry}

이번 단원의 핵심은 다음 문장으로 정리할 수 있다.

불리언 마스크는 단순한 참·거짓 배열이 아니라, Pandas에서는 Index와 dtype을 갖고 데이터 선택 및 변환에 참여하는 객체일 수 있다.

\subsubsection{불리언 마스크의 객체 모델}

먼저 가장 기본적인 부릴언 마스크를 생성해 보자.

\begin{dhknote}
\begin{dhksmry}[centered]
독립 실행 코드 1 - 비교식의 반환 객체
\end{dhksmry}
\end{dhknote}

\begin{dhkcode}
import pandas as pd

df = pd.DataFrame(
    {
        "score": [90, 70, 85],
    },
    index=pd.Index(
        ["A", "B", "C"],
        name="student_id",
    ),
)

mask = df["score"] >= 80

print(mask)
print(type(mask).__name__)
print(mask.dtype)
print(mask.index.equals(df.index))
\end{dhkcode}

예상 출력:

\begin{dhkzsh}
student_id
A     True
B    False
C     True
Name: score, dtype: bool

Series
bool
True
\end{dhkzsh}

이제 객체 생성 과정과 반환 과정을 분석해 보자.

\begin{dhknote}
\begin{dhksmry}[centered]
1. df["score"]의 평가
\end{dhksmry}
\end{dhknote}

\begin{dhkcode}
df["score"]
\end{dhkcode}

DataFrame의 \zsh{__getitem__} 메커니즘을 통해 \zsh{"score"} 열을 선택한다.

그 결과 Series 객체를 얻는다.

\begin{dhknote}
\begin{dhksmry}[centered]
2. >= 80의 평가
\end{dhksmry}
\end{dhknote}

\begin{dhkcode}
df["score"] >= 80
\end{dhkcode}

Python의 비교 연산자 \zsh{>=}에 대응하는 Pandas의 연산 구현이 호출된다.

Pandas는 Series의 각 원소와 스칼라 80을 비교하여 결과를 계산한다.

여기서 80은 Index를 갖는 객체가 아니므로, Series의 각 원소에 같은 기준값을 적용한다.

\begin{dhknote}
\begin{dhksmry}[centered]
3. 새로운 Series 반환
\end{dhksmry}
\end{dhknote}

비교 결과는 단일 불리언 값이 아니다.

\begin{dhkzsh}
A → True
B → False
C → True
\end{dhkzsh}

이 결과를 관리하는 Series가 생성되고 \zsh{mask}라는 이름에 바인딩된다.

원본 Series의 Index 레이블도 결과에 유지된다.

\begin{dhknote}
\begin{dhksmry}[centered]
4. 마스크와 원본 DataFrame의 관계
\end{dhksmry}
\end{dhknote}

\zsh{mask}는 \zsh{df}와 동일한 객체가 아니다.

또한 \zsh{mask}는 원본 점수들을 저장하는 것이 아니라, 비교 연산의 결과를 저장한다.

따라서 원본의 점수를 나중에 변경한다고 해서 이미 계산된 \zsh{mask}가 자동으로 다시 계산되는 것은 아니다.

\begin{dhkcode}
df.loc["B", "score"] = 95

print(mask.loc["B"])
print((df["score"] >= 80).loc["B"])
\end{dhkcode}

예상 출력:

\begin{dhkzsh}
False
True
\end{dhkzsh}

첫 번째 출력은 과거에 계산한 \zsh{mask}의 값이다.

두 번째 출력은 변경된 DataFrame을 대상으로 비교식을 다시 평가한 결과이다.

즉, 불리언 마스크는 조건식을 기억하고 자동 갱신하는 객체가 아니라, 특정 시점에 계산된 조건의 결과를 담은 객체다.

이 차이는 Python의 함수 호출 결과와 함수 객체를 구분했던 내용과도 연결된다.

\subsubsection{여러 조건을 결합하는 방법 - \zsh{&, |, ~}}

실제 데이터 분석에서는 하나의 조건만 사용하는 경우보다 여러 조건을 결합하는 경우가 많다.

예를 들어 문서 분류 데이터셋에서 다음 조건을 만족하는 문서를 찾는다고 가정하자.

\begin{dhksmry}
    - 데이터 분할이 "valid"일 것

    - 토큰 길이가 300을 초과할 것

    - 언어가 "ko"일 것
\end{dhksmry}

Python에서는 일반적으로 \zsh{and}, \zsh{or}, \zsh{not}을 사용해 논리 연산을 수행한다.

하지만 Pandas Series에 대해서는 원소별 논리 연산을 수행하기 위해 \zsh{&}, \zsh{|}, \zsh{~}를 사용해야 한다.

\begin{dhknote}
\begin{dhksmry}[centered]
독립 실행 코드 2 - 공통 분석 데이터셋
\end{dhksmry}
\end{dhknote}

이후 예제에서 사용할 DataFrame을 먼저 정의하자.

\begin{dhkcode}
import pandas as pd

df = pd.DataFrame(
    {
        "language": ["ko", "en", "ko", "en", "ko"],
        "label": [1, 0, 1, None, 0],
        "prediction": [1, 1, 0, 1, None],
        "token_length": [120, 420, None, 300, 510],
        "split": ["train", "valid", "valid", "train", "valid"],
    },
    index=pd.Index(
        ["d1", "d2", "d3", "d4", "d5"],
        name="doc_id",
    ),
).astype({
    "label": "Int64",
    "prediction": "Int64",
    "token_length": "Int64",
})

print(df)
\end{dhkcode}

예상 출력:

\begin{dhkzsh}
       language  label  prediction  token_length  split
doc_id
d1           ko      1           1           120  train
d2           en      0           1           420  valid
d3           ko      1           0          <NA>  valid
d4           en   <NA>           1           300  train
d5           ko      0        <NA>           510  valid
\end{dhkzsh}

이제 여러 조건을 결합해 보자.

\begin{dhkcode}
valid = df["split"] == "valid"
korean = df["language"] == "ko"
long = df["token_length"] > 300

mask = valid & korean & long

print(mask.tolist())
print(df.loc[mask].index.tolist())
\end{dhkcode}

예상 출력:

\begin{dhkzsh}
[False, False, <NA>, False, True]
['d5']
\end{dhkzsh}

세 조건은 각각 Series를 반환한다.

\zsh{&} 연산자는 각 Series의 원소별 논리 AND 연산을 수행해한다.

여기서 중요한 것은 다음 두 표현식의 차이다.

\begin{dhkcode}
# Pandas Series의 원소별 논리 연산
mask = valid & korean

# Python 객체 자체의 진리값 평가
mask = valid and korean
\end{dhkcode}

두 번째 표현식은 일반적인 다중 원소 Series에 대해 다음 오류를 발생시킨다.

\begin{dhkzsh}
ValueError: The truth value of a Series is ambiguous.
\end{dhkzsh}

왜 이런 오류가 발생할까?

Python의 \zsh{and}는 피연산자 객체의 진리값을 평가하여 연산한다.

그런데 Series에는 여러 개의 불리언 값이 존재하므로 Series 객체 전체를 참으로 볼지, 거짓으로 볼지 결정할 수 없다.

반면 Pandas Series의 \zsh{&} 연산은 원소별 논리 연산을 수행하도록 구현되어 있다.

\begin{dhktable}
|연산자|	의미|	Pandas Series에서의 사용|
|:---|:---|:---|
|&|	원소별 AND|	여러 조건을 모두 만족|
|│|	원소별 OR	|하나 이상의 조건을 만족|
|~|	원소별 NOT	|조건 결과를 반전|
|and|	Python 논리 AND|	일반적인 Series끼리 직접 사용 불가|
|or|	Python 논리 OR	|일반적인 Series끼리 직접 사용 불가|
|not|	Python 논리 NOT	|일반적인 Series에 직접 사용 불가|
\end{dhktable}

\begin{dhknote}
\begin{dhksmry}[centered]
비교식에 괄호가 필요한 이유
\end{dhksmry}
\end{dhknote}

다음 표현식을 살펴보자.

\begin{dhkcode}
mask = (
    (df["token_length"] > 300)
    & (df["label"] == 0)
)
\end{dhkcode}

이것은 올바른 표현식이다.

하지만 다음처럼 작성하면 문제가 발생할 수 있다.

\begin{dhkcode}
mask = df["token_length"] > 300 & df["label"] == 0
\end{dhkcode}

Python에서는 \zsh{&}의 연산자 우선순위가 비교 연산자보다 높다.

따라서 개발자가 의도한 다음 구조로 평가되지 않는다.

\begin{dhkcode}
(df["token_length"] > 300) & (df["label"] == 0)
\end{dhkcode}

실제로는 \zsh{300 & df["label"]}과 같은 부분이 먼저 평가되는 등 의도하지 않은 연산 구조가 만들어진다.

그 결과 잘못된 비교나 \zsh{ValueError}가 발생할 수 있다.

따라서 Pandas의 복합 조건에서는 각 비교식을 괄호로 묶는 습관이 중요하다.

또한 \zsh{&}와 \zsh{|}는 Python의 \zsh{and}, \zsh{or}처럼 피연산자를 단락 평가(short-circuit evaluation)하지 않는다.

양쪽 조건식을 평가한 뒤 원소별 연산을 수행하므로, 한쪽 조건이 거짓이라는 이유로 다른 쪽의 오류 가능성이 사라지는 것은 아니다.

\subsubsection{Nullable Boolean}과 결측갑스이 선택 규칙

\zsh{D.3}에서는 Nullable Boolean이 \zsh{True}, \zsh{False}, \zsh{<NA>}를 표현할 수 있다는 사실을 배웠다.

그렇다면 \zsh{<NA>}가 포함된 마스크를 \zsh{.loc}에 전달하면 어떻게 될까?

\begin{dhknote}
\begin{dhksmry}[centered]
독립 실행 코드 3 - 결측을 포함한 마스크
\end{dhksmry}
\end{dhknote}

앞서 정의한 \zsh{df}를 그대로 사용하자.

\begin{dhkcode}
mask = df["token_length"] > 300

print(mask.dtype)
print(mask.tolist())

print(df.loc[mask].index.tolist())
print(df.loc[~mask].index.tolist())
print(df.loc[mask.isna()].index.tolist())
\end{dhkcode}

예상 출력:

\begin{dhkzsh}
boolean
[False, True, <NA>, False, True]
['d2', 'd5']
['d1', 'd4']
['d3']
\end{dhkzsh}

이 결과를 주의 깊게 살펴봐야 한다.

\zsh{d3}의 토큰 길이는 결측이다.

따라서 다음 비교식의 결과도 \zsh{<NA>}가 된다.

\begin{dhkcode}
df["token_length"] > 300
\end{dhkcode}

Pandas에서 Nullable Boolean 마스크를 행 선택에 사용할 때는 \zsh{<NA>}가 있는 위치를 선택하지 않는다.

따라서 \zsh{df.loc[mask]}는 \zsh{d2}, \zsh{d5}만 반환한다.

그런데 다음 표현식은 어떨까?

\begin{dhkcode}
df.loc[~mask]
\end{dhkcode}

\zsh{~}는 \zsh{True}와 \zsh{False}를 반전하지만 \zsh{<NA>}는 그대로 \zsh{<NA>}로 유지한다.

따라서 \zsh{d3}는 반대 조건으로 선택할 때도 제외된다.

이것은 매우 중요한 사실이다.

결측이 포함된 마스크에서는 \zsh{mask}로 선택한 행과 \zsh{~mask}로 선택한 행의 합집합이 원본의 모든 행이 아닐 수 있다.

\begin{dhkgraph}[uniform]
node n1 (1,1) [accent] :: True // d2, d5 // 조건 만족
node n2 (1,2) [blue] :: False // d1, d4 // 조건 불만족
node n3 (1,3) [orange] :: <NA> // d3 // 판정 불가
explain :: 조건 결과에 따른 세 가지 상태
\end{dhkgraph}

따라서 다음 세 집합을 구분해야 한다.

\begin{dhkcode}
passed = df.loc[mask]
failed = df.loc[~mask]
unknown = df.loc[mask.isna()]
\end{dhkcode}

이 방식은 조건을 만족하는 데이터와 만족하지 않는 데이터뿐 아니라, 판정 자체가 불가능한 데이터도 별도로 보존한다.

\begin{dhknote}
\begin{dhksmry}[centered]
fillna()로 결측을 처리하면?
\end{dhksmry}
\end{dhknote}

\begin{dhkcode}
print(mask.fillna(False).tolist())
print(mask.fillna(True).tolist())
\end{dhkcode}

예상 출력:

\begin{dhkzsh}
[False, True, False, False, True]
[False, True, True, False, True]
\end{dhkzsh}

\zsh{fillna(False)}는 판정할 수 없는 데이터를 조건 불만족으로 처리한다.

반면 \zsh{fillna(True)}는 판정할 수 없는 데이터를 조건 만족으로 처리한다.

어느 쪽이 올바른지는 분석 목적에 따라 다르다.

예를 들어 토큰 길이가 300을 초과하는 문서를 제외하는 전처리라면, 길이를 알 수 없는 문서를 함께 제외할지 별도로 검수할지 결정해야 한다.

중요한 것은 결측값을 \zsh{False}로 바꾸는 행위가 단순한 기술적 변환이 아니라 데이터 처리 정책이라는 점이다.

\subsubsection{\zsh{isin()} - 하나의 값이 여러 후보에 속하는가?}

지금까지는 하나의 값과 비교했다.

\begin{dhkcode}
df["language"] == "ko"
\end{dhkcode}

그런데 분석 대상 언어가 \zsh{"ko"} 또는 \zsh{"en"}처럼 여러 개라면 어떻게 할까?

이때 사용할 수 있는 메서드가 \zsh{isin()}이다.

\begin{dhknote}
\begin{dhksmry}[centered]
독립 실행 코드 4 - 여러 값에 대한 소속 검사
\end{dhksmry}
\end{dhknote}

\begin{dhkcode}
import pandas as pd

s = pd.Series(
    ["ko", "en", "ja", None],
    dtype="string[python]",
)

mask = s.isin(["ko", "en"])

print(mask.tolist())
print(s.loc[mask].tolist())
\end{dhkcode}

예상 출력:

\begin{dhkzsh}
[True, True, False, False]
['ko', 'en']
\end{dhkzsh}

\zsh{isin()}은 각 원소가 전달된 후보 집합에 포함되는지 검사한다.

따라서 다음과 같은 의미를 갖는다.

\begin{dhkcode}
s.isin(["ko", "en"])
\end{dhkcode}

각 원소가 \zsh{"ko"} 또는 \zsh{"en"}에 해당하는가?

일반적인 단일 비교와 비교하면 다음과 같다.

\begin{dhktable}
|표현식|	의미|
|:---|:---|
|s == "ko"|	"ko"와 같은가?|
|s.isin(["ko", "en"])|	"ko" 또는 "en" 중 하나인가?|
|~s.isin(["ko", "en"])|	두 후보 중 어느 것에도 속하지 않는가?|
\end{dhktable}

여기서 주의할 점이 있다.

\zsh{isin()}은 문자열 내부의 부분 문자열을 찾는 메서드가 아니다.

예를 들어 \zsh{"korean"}이라는 값이 있을 때 \zsh{"ko"}가 부분 문자열로 포함되어 있다는 이유만으로 \zsh{isin(["ko"])}가 참이 되지는 않는다.

부분 문자열 검색은 이후 \zsh{D.10}에서 다룰 \zsh{.str.contains()}와 같은 문자열 접근자를 사용해야 한다.

\begin{dhknote}
\begin{dhksmry}[centered]
Nullable 자료형에서 isin()과 비교식의 차이
\end{dhksmry}
\end{dhknote}

\begin{dhkcode}
s = pd.Series(
    [1, None, 0],
    dtype="Int64",
)

print((s == 1).tolist())
print(s.isin([1]).tolist())
print(s.isin([pd.NA]).tolist())
\end{dhkcode}

예상 출력:

\begin{dhkzsh}
[True, <NA>, False]
[True, False, False]
[False, True, False]
\end{dhkzsh}

여기서 중요한 차이가 있다.

\zsh{==}는 Nullable 정수형의 비교 의미론을 사용하므로 결측 위치에 \zsh{<NA>}를 반환한다.

반면 \zsh{isin()}은 후보 집합에 포함되는지를 판단한다.

이 예제에서는 후보에 결측이 없으면 결측 위치가 \zsh{False}이고, \zsh{pd.NA}를 후보로 전달하면 해당 결측 위치가 \zsh{True}가 된다.

따라서 \zsh{isin()}을 단순히 여러 개의 \zsh{==} 비교를 연결한 것과 모든 경우에 동등하다고 가정하면 안 된다.

\subsubsection{\zsh{where()}와 \zsh{mask()} - 선택과 값 변경의 차이}

이제 이번 학습의 중요한 구분을 살펴보자.

\begin{dhkcode}
df.loc[condition]
\end{dhkcode}

이 표현식은 조건을 만족하는 행을 선택한다.

그런데 다음 표현식은 다르게 작동한다.

\begin{dhkcode}
df.where(condition)
\end{dhkcode}

\zsh{where()}는 일반적으로 원본 객체의 형태를 유지하면서, 조건을 만족하지 않는 위치의 값을 다른 값으로 대체하는 메서드다.

반대로 \zsh{mask()}는 조건을 만족하는 위치의 값을 대체한다.

\begin{dhknote}
\begin{dhksmry}[centered]
독립 실행 코드 5 - where()와 mask()
\end{dhksmry}
\end{dhknote}

앞서 정의한 DataFrame에서 토큰 길이 Series를 선택한다.

\begin{dhkcode}
s = df["token_length"]

condition = s <= 300

print(condition.tolist())

print(s.where(condition, other=0).tolist())
print(s.mask(condition, other=0).tolist())
\end{dhkcode}

예상 출력:

\begin{dhkzsh}
[True, False, <NA>, True, False]
[120, 0, 0, 300, 0]
[0, 420, 0, 0, 510]
\end{dhkzsh}

각 원소에서 무슨 일이 발생했는지 분석해 보자.

\begin{dhktable}
|문서|	원래 값|	조건 <=300|	where(..., 0)|	mask(..., 0)|
|:---|:---|:---|:---|:---|
|d1|	120|	True|	120|	0|
|d2|	420|	False|	0|	420|
|d3|	<NA>|	<NA>|	0|	0|
|d4|	300|	True|	300|	0|
|d5|	510|	False|	0|	510|
\end{dhktable}

\zsh{where()}는 조건이 참인 위치의 원래 값을 유지하고, 나머지를 대체한다.

반면 \zsh{mask()}는 조건이 참인 위치를 대체하고, 거짓인 위치의 값을 유지한다.

여기서 결측 조건의 처리가 중요하다.

Pandas의 \zsh{where()}는 조건이 \zsh{<NA>}인 위치를 대체 대상으로 처리한다.

\zsh{mask()} 역시 조건이 \zsh{<NA>}인 위치를 대체 대상으로 처리한다.

따라서 두 경우 모두 \zsh{d3}의 값이 0으로 바뀌었다.

\begin{dhknote}
\begin{dhksmry}[centered]
where()와 mask()는 항상 완전한 반대일까?
\end{dhksmry}
\end{dhknote}

일반적인 참·거짓 조건에 대해서는 다음과 같은 관계로 이해할 수 있다.

\begin{dhkcode}
s.mask(condition, other=0)
s.where(~condition, other=0)
\end{dhkcode}

하지만 Nullable Boolean의 \zsh{<NA>}가 포함되면 조건 반전 후에도 \zsh{<NA>}가 남는다.

이때 Pandas의 결측 조건 처리 규칙까지 함께 고려해야 하므로, 단순한 2값 논리의 반대 연산으로만 이해해서는 안 된다.

\begin{dhknote}
\begin{dhksmry}[centered]
행 필터링과의 차이
\end{dhksmry}
\end{dhknote}

다음 세 표현식을 비교해 보자.

\begin{dhkcode}
selected = s.loc[condition]
replaced_false = s.where(condition, other=0)
replaced_true = s.mask(condition, other=0)
\end{dhkcode}

\begin{dhktable}
|표현식|	결과의 길이|	동작|
|:---|:---|:---|
|s.loc[condition]|	조건에 따라 감소|	조건을 만족하는 원소만 선택|
|s.where(condition, 0)|	원본과 동일|	조건 불만족 위치를 대체|
|s.mask(condition, 0)|	원본과 동일|	조건 만족 위치를 대체|
\end{dhktable}

즉, \zsh{loc}는 데이터를 선택하는 연산이고, \zsh{where}와 \zsh{mask}는 값을 조건부로 유지하거나 대체하는 연산이다.

\begin{dhknote}
\begin{dhksmry}[centered]
결측을 유지하고 싶다면?
\end{dhksmry}
\end{dhknote}

\begin{dhkcode}
print(s.where(condition).tolist())
print(s.mask(condition).tolist())
\end{dhkcode}

예상 출력:

\begin{dhkzsh}
[120, <NA>, <NA>, 300, <NA>]
[<NA>, 420, <NA>, <NA>, 510]
\end{dhkzsh}

\zsh{other}를 생략하면 일반적으로 조건에 따라 대체되는 위치에 결측값이 사용된다.

이 경우에는 원본 dtype과 결측 표현 방식에 따라 결과 dtype이 달라질 수 있다.

이번 예제의 \zsh{s}는 Nullable 정수형 \zsh{Int64}이므로 \zsh{<NA>}를 표현할 수 있다.

또한 \zsh{where()}와 \zsh{mask()}는 기본적으로 결과 객체를 반환한다.

반환값을 별도의 이름에 바인딩하는 것만으로 원본 Series가 변경되는 것은 아니다.

\subsubsection{\zsh{query()} - 문자열로 조건식을 표현하기}

Pandas는 조건을 문자열로 작성할 수 있는 \zsh{DataFrame.query()} 메서드도 제공한다.

\begin{dhknote}
\begin{dhksmry}[centered]
독립 실행 코드 6 - query() 사용
\end{dhksmry}
\end{dhknote}

\begin{dhkcode}
threshold = 400

result = df.query(
    "token_length > @threshold and language == 'ko'",
    engine="python",
)

print(result.index.tolist())
\end{dhkcode}

예상 출력:

\begin{dhkzsh}
['d5']
\end{dhkzsh}

이 코드는 다음 조건을 표현한다.

\begin{dhksmry}
    - token_length가 외부 변수 threshold보다 클 것

    - language가 "ko"일 것
\end{dhksmry}

여기서 \zsh{@threshold}는 \zsh{query()} 표현식 외부에 존재하는 Python 변수를 참조하기 위한 문법이다.

또한 \zsh{query()} 내부에서는 \zsh{and}, \zsh{or}, \zsh{not}과 같은 표현을 사용할 수 있다.

이는 일반적인 Series 조건식을 Python 코드로 직접 작성할 때 \zsh{&}, \zsh{|}, \zsh{~}를 사용해야 하는 것과 다르다.

\begin{dhknote}
\begin{dhksmry}[centered]
일반적인 불리언 마스크와 비교
\end{dhksmry}
\end{dhknote}

다음 두 표현식은 현재 예제에서 같은 행을 선택한다.

\begin{dhkcode}
result1 = df.loc[
    (df["token_length"] > threshold)
    & (df["language"] == "ko")
]

result2 = df.query(
    "token_length > @threshold and language == 'ko'",
    engine="python",
)
\end{dhkcode}

하지만 실행 메커니즘은 같지 않다.

첫 번째 표현식은 Python에서 Series 비교 연산을 수행하고, 그 결과를 \zsh{.loc}에 전달한다.

두 번째 표현식은 \zsh{query()}가 문자열 표현식을 해석하여 조건을 평가하고, 그 조건에 맞는 행을 선택한다.

따라서 \zsh{query()}는 단순히 \zsh{.loc}의 다른 표기법이 아니라 표현식 평가 기능을 포함하는 메서드다.

\begin{dhknote}
\begin{dhksmry}[centered]
열 이름에 공백이 있다면?
\end{dhksmry}
\end{dhknote}

\begin{dhkcode}
sample = pd.DataFrame({
    "token length": [100, 500],
})

result = sample.query(
    "`token length` > 300"
)

print(result.index.tolist())
\end{dhkcode}

예상 출력:

\begin{dhkzsh}
[1]
\end{dhkzsh}

열 이름에 공백이 있는 경우에는 백틱(backtick)을 사용해 참조할 수 있다.

다만 모든 특수한 열 이름을 제한 없이 지원하는 것은 아니다.

\begin{dhknote}
\begin{dhksmry}[centered]
SQL의 WHERE와 같은가?
\end{dhksmry}
\end{dhknote}

문법적으로 비슷한 부분은 있지만 동일한 시스템은 아니다.

\begin{dhktable}
|구분|	Pandas query()|	SQL WHERE|
|:---|:---|:---|
|대상|	메모리의 DataFrame|	데이터베이스의 테이블·쿼리 결과|
|조건 표현|	Pandas 표현식 문법|	SQL 조건식|
|외부 변수|	@variable|	DB 드라이버의 매개변수 바인딩 등|
|결측 의미론|	Pandas dtype과 연산 규칙|	SQL의 NULL 및 3값 논리|
|실행|	Pandas 표현식 평가 엔진|	데이터베이스 쿼리 엔진|
\end{dhktable}

특히 \zsh{query()}에 전달하는 문자열은 평가 대상 표현식이므로, 신뢰할 수 없는 외부 입력을 그대로 삽입해서는 안 된다.

또한 \zsh{engine="python"}은 Python 평가 엔진을 사용하도록 명시한 것이다.

Nullable dtype을 사용하는 예제에서 동작을 명확하게 확인하기 위해 지정했지만, 데이터 크기와 표현식에 따라 다른 엔진보다 느릴 수 있다.

\subsubsection{불리언 마스크의 Index alignment}

\zsh{D.5}에서는 두 Series의 산술 연산에서 Index alignment가 발생한다는 사실을 배웠다.

불리언 마스크를 \zsh{.loc}에 전달할 때도 Index 레이블이 중요하다.

\begin{dhknote}
\begin{dhksmry}[centered]
독립 실행 코드 7 - 같은 값, 다른 선택 결과
\end{dhksmry}
\end{dhknote}

\begin{dhkcode}
import pandas as pd

df = pd.DataFrame(
    {"score": [10, 20, 30, 40, 50]},
    index=["d1", "d2", "d3", "d4", "d5"],
)

mask = pd.Series(
    [True, False, True, False, True],
    index=["d5", "d1", "d2", "d4", "d3"],
)

by_label = df.loc[mask]

by_position = df.iloc[
    mask.to_numpy()
]

print(by_label.index.tolist())
print(by_position.index.tolist())
\end{dhkcode}

예상 출력:

\begin{dhkzsh}
['d2', 'd3', 'd5']
['d1', 'd3', 'd5']
\end{dhkzsh}

두 결과가 달라졌다.

왜 그럴까?

\zsh{.loc[mask]}는 마스크의 Index를 DataFrame의 Index에 대응시킨다.

따라서 마스크에서 \zsh{True}인 레이블 \zsh{d5}, \zsh{d2}, \zsh{d3}에 해당하는 행을 선택한다.

반환된 행은 원본 DataFrame의 순서에 따라 \zsh{d2}, \zsh{d3}, \zsh{d5}로 나타난다.

반면 \zsh{mask.to_numpy()}는 Index 레이블 정보를 제거하고 불리언 값만 배열로 반환한다.

\begin{dhkzsh}
[True, False, True, False, True]
\end{dhkzsh}

따라서 \zsh{.iloc}에서는 첫 번째, 세 번째, 다섯 번째 위치를 선택하게 된다.

\begin{dhknote}
\begin{dhksmry}[centered]
마스크의 레이블이 부족하다면?
\end{dhksmry}
\end{dhknote}

\begin{dhkcode}
bad_mask = pd.Series(
    [True, False],
    index=["d1", "d2"],
)

df.loc[bad_mask]
\end{dhkcode}

이 코드는 다음 오류를 발생시킨다.

\begin{dhkzsh}
IndexingError: Unalignable boolean Series provided as indexer ...
\end{dhkzsh}

DataFrame의 모든 행에 대응하는 불리언 값을 구성할 수 없기 때문이다.

여기서 주의할 점은 \zsh{.loc}의 불리언 마스크 처리와 \zsh{where()}의 조건 정렬이 완전히 같은 정책을 갖는 것은 아니라는 사실이다.

예를 들어 \zsh{where()}는 대응되지 않는 조건 위치를 대체 대상으로 처리할 수 있다.

반면 \zsh{.loc}은 필요한 레이블에 대응하는 불리언 값이 없으면 오류를 발생시킬 수 있다.

따라서 Pandas의 모든 메서드가 Index alignment를 수행한다고 해서, 정합하지 않은 레이블을 처리하는 방식까지 동일한 것은 아니다.

\subsubsection{조건부 선택과 조건부 할당은 다르다}

이전 학습에서는 데이터 조회와 상태 변경을 구분했다.

불리언 마스크에서도 이 차이를 유지해야 한다.

\begin{dhknote}
\begin{dhksmry}[centered]
독립 실행 코드 8 - 선택과 상태 변경
\end{dhksmry}
\end{dhknote}

\begin{dhkcode}
import pandas as pd

df = pd.DataFrame(
    {
        "score": [90, 70, 85],
        "status": ["pending", "pending", "pending"],
    },
    index=["A", "B", "C"],
)

mask = df["score"] >= 80

selected = df.loc[mask]

df.loc[mask, "status"] = "passed"

print(selected["status"].tolist())
print(df["status"].tolist())
\end{dhkcode}

Pandas 2.2.3에서 Copy-on-Write를 활성화한 환경 또는 Pandas 3.x의 Copy-on-Write 동작을 기준으로 보면 결과는 다음과 같다.

\begin{dhkzsh}
['pending', 'pending']
['passed', 'pending', 'passed']
\end{dhkzsh}

여기서 객체 관계를 분석해 보자.

먼저 다음 표현식은 조건을 만족하는 행을 선택해 결과 DataFrame을 반환한다.

\begin{dhkcode}
selected = df.loc[mask]
\end{dhkcode}

반면 다음 표현식은 원본 DataFrame의 지정된 위치에 값을 할당한다.

\begin{dhkcode}
df.loc[mask, "status"] = "passed"
\end{dhkcode}

두 표현식의 역할은 다르다.

또한 \zsh{selected}가 원본 \zsh{df}와 내부 데이터를 공유하는지, 독립적으로 보유하는지는 Pandas 버전과 Copy-on-Write 설정 등에 영향을 받을 수 있다.

따라서 \zsh{selected}가 원본의 변경을 자동으로 반영한다고 가정해서는 안 된다.

실무에서 명확한 의도가 중요하다면 다음처럼 구분하는 편이 좋다.

\begin{dhkcode}
# 조건에 맞는 데이터 추출
selected = df.loc[mask].copy()

# 원본 DataFrame의 특정 위치 변경
df.loc[mask, "status"] = "passed"
\end{dhkcode}

이 부분은 \zsh{D.7}에서 객체 복사, 내부 데이터 공유, Copy-on-Write와 함께 더 자세히 다룰 것이다.

\subsubsection{Transformer 응용 - 오류 샘플과 길이 초과 문서 추출}

이제 이번 단원의 개념을 하나의 분석 과정에 통합해 보자.

문서 분류 모델의 평가 결과를 분석한다고 가정하자.

분석 목표는 두 가지다.

첫째, 검증 데이터 중 정답 레이브로가 예측 레이블이 다른 문서를 찾는 것이다.

둘째, 토큰 길이가 300을 초과하는 문서를 찾되, 길이를 알 수 없는 문서를 별도로 보존하는 것이다.

\zsh{D.6.4}에서 정의한 공통 DataFrame을 사용하자.

\begin{dhknote}
\begin{dhksmry}[centered]
독립 실행 코드 9 - 오류 문서 분석
\end{dhksmry}
\end{dhknote}

\begin{dhkcode}
import pandas as pd

df = pd.DataFrame(
    {
        "language": ["ko", "en", "ko", "en", "ko"],
        "label": [1, 0, 1, None, 0],
        "prediction": [1, 1, 0, 1, None],
        "token_length": [120, 420, None, 300, 510],
        "split": ["train", "valid", "valid", "train", "valid"],
    },
    index=pd.Index(
        ["d1", "d2", "d3", "d4", "d5"],
        name="doc_id",
    ),
).astype({
    "label": "Int64",
    "prediction": "Int64",
    "token_length": "Int64",
})

valid = df["split"].eq("valid")

known = (
    df["label"].notna()
    & df["prediction"].notna()
)

wrong = df["label"].ne(
    df["prediction"]
)

error_mask = valid & known & wrong

errors = df.loc[
    error_mask,
    ["language", "label", "prediction"],
]

print(errors)
\end{dhkcode}

예상 출력:

\begin{dhkzsh}
       language  label  prediction
doc_id
d2           en      0           1
d3           ko      1           0
\end{dhkzsh}

\begin{dhknote}
\begin{dhksmry}[centered]
실행 과정 분석
\end{dhksmry}
\end{dhknote}

\begin{dhknote}
\begin{dhksmry}[centered]
1단계 - 검증 데이터 판정
\end{dhksmry}
\end{dhknote}

\begin{dhkcode}
valid = df["split"].eq("valid")
\end{dhkcode}

검증 데이터에 해당하는 행을 나타내는 불리언 Series를 생성한다.

\begin{dhknote}
\begin{dhksmry}[centered]
2단계 - 평가 가능한 문서 판정
\end{dhksmry}
\end{dhknote}

\begin{dhkcode}
known = (
    df["label"].notna()
    & df["prediction"].notna()
)
\end{dhkcode}

정답과 예측이 모두 존재하는 문서만 평가 대상으로 인정한다.

여기서 결측을 무조건 오답으로 처리하지 않았다는 점이 중요하다.

\begin{dhknote}
\begin{dhksmry}[centered]
3단계 - 예측 오류 판정
\end{dhksmry}
\end{dhknote}

\begin{dhkcode}
wrong = df["label"].ne(df["prediction"])
\end{dhkcode}

\zsh{ne()}는 \zsh{not equal} 비교 메서드다.

두 Series의 값을 비교해 서로 다르면 \zsh{True}를 반환한다.

Nullable 정수형이므로 어느 한쪽이 결측인 경우에는 \zsh{<NA>}가 발생할 수 있다.

\begin{dhknote}
\begin{dhksmry}[centered]
4단계 - 최종 조건 결합
\end{dhksmry}
\end{dhknote}

\begin{dhkcode}
error_mask = valid & known & wrong
\end{dhkcode}

세 조건을 결합해 검증 데이터이면서, 정답과 예측이 존재하고, 두 값이 서로 다른 문서만 선택한다.

이때 \zsh{known}이 \zsh{False}인 행은 Nullable Boolean의 AND 연산 규칙에 따라 최종 결과가 \zsh{False}로 확정될 수 있다.

\begin{dhknote}
\begin{dhksmry}[centered]
5단계 - 오류 샘플 추출
\end{dhksmry}
\end{dhknote}

\begin{dhkcode}
errors = df.loc[
    error_mask,
    ["language", "label", "prediction"],
]
\end{dhkcode}

마지막으로 \zsh{.loc}을 사용해 조건을 만족하는 행과 필요한 열만 선택한다.

결과는 두 행을 가진 DataFrame이다.

원본 문서 ID인 \zsh{d2}, \zsh{d3}도 Index로 유지된다.

이렇게 하면 모델의 오류를 분석할 때 해당 문서의 원본 텍스트나 메타데이터를 다시 추적할 수 있다.

\begin{dhknote}
\begin{dhksmry}[centered]
길이 초과 문서와 길이 미확정 문서 구분
\end{dhksmry}
\end{dhknote}

이번에는 토큰 길이를 분석해 보자.

\begin{dhkcode}
long_mask = df["token_length"] > 300

long_docs = df.loc[long_mask]
unknown_docs = df.loc[long_mask.isna()]

print(long_docs.index.tolist())
print(unknown_docs.index.tolist())

print(long_mask.notna().sum())
print(len(df))
\end{dhkcode}

예상 출력:

\begin{dhkzsh}
['d2', 'd5']
['d3']
4
5
\end{dhkzsh}

총 다섯 문서 중 네 문서는 토큰 길이 조건을 판정할 수 있다.

그중 두 문서의 길이가 300을 초과하고, 한 문서는 길이를 알 수 없다.

여기서 단순히 다음과 같이 계산하면 문제가 생길 수 있다.

\begin{dhkcode}
len(long_docs) / len(df)
\end{dhkcode}

결과는 \zsh{2 / 5 = 0.4}다.

하지만 이것은 전체 문서 중 길이 초과로 확인된 문서의 비율이지, 길이가 관측된 문서 중 길이 초과 비율과는 다르다.

후자를 계산하려면 다음과 같이 해야 한다.

\begin{dhkcode}
long_mask.sum() / long_mask.notna().sum()
\end{dhkcode}

결과는 \zsh{2 / 4 = 0.5}다.

두 지표는 분모의 정의가 다르다.

\begin{dhktable}
|지표	|계산|	값|
|:---|:---|:---|
|전체 문서 중 길이 초과 확인 비율|	2 / 5	|40 퍼센트|
|길이가 확인된 문서 중 초과 비율|	2 / 4	|50 퍼센트|
|토큰 길이 미확정 비율|	1 / 5	|20 퍼센트|
\end{dhktable}

이것은 단순한 코드 차이가 아니라 분석 대상 모집단과 관측 가능한 표본을 어떻게 정의하는가의 문제다.

따라서 결측값을 포함한 불리언 마스크에서는 참인 원소의 개수뿐 아니라, 판정 가능한 원소의 개수도 함께 관리해야 한다.

\subsubsection{주요 메서드 비교}

이번 학습에서 배운 인터페이스를 비교해 보자.

\begin{dhktable}
|표현식	|주요 목적	|반환 객체|
|:---|:---|:---|
|s > value|	비교 조건 생성|	불리언 Series|
|a & b|	원소별 AND|	불리언 Series|
|a │ b|	원소별 OR|	불리언 Series|
|~a	|원소별 NOT|	불리언 Series|
|s.isin(values)	|후보 집합 소속 검사|	불리언 Series|
|df.loc[mask]|	조건에 맞는 행 선택	|DataFrame|
|s.where(cond, other)|	조건이 거짓인 위치 대체|	Series|
|s.mask(cond, other)|	조건이 참인 위치 대체|	Series|
|df.query(expr)	|문자열 조건식으로 행 선택|	DataFrame|
|df.loc[mask, col] = value|	조건부 할당	|기존 DataFrame 상태 변경|
\end{dhktable}

여기서 세 가지를 구분하면 대부분의 조건부 분석 코드를 설명할 수 있다.

조건 생성은 비교식과 논리 연산을 통해 불리언 객체를 만드는 과정이다.

조건부 선택은 불리언 객체를 사용해 필요한 행이나 원소를 추출하는 과정이다.

조건부 변환은 원본의 형태를 유지하면서 조건에 따라 값을 대체하거나, 특정 위치의 상태를 변경하는 과정이다.

이 세 과정은 함께 사용되지만 동일한 연산은 아니다.

\subsubsection{자주 발생하는 오해와 경계 조건}

\begin{dhktable}
|오해	|정확한 이해|
|:---|:---|
|불리언 마스크는 항상 NumPy 배열이다|	Pandas에서는 Index와 dtype을 가진 Series일 수 있다|
|and, or로 Series 조건을 결합하면 된다|	일반적인 원소별 조건에는 &, │를 사용해야 한다
|~mask는 mask로 선택되지 않은 모든 행을 선택한다|	<NA>는 반전 후에도 결측이므로 양쪽에서 제외될 수 있다|
|isin()은 여러 개의 ==를 단순 결합한 것이다|	결측값 처리 등에서 차이가 있다
|where()는 조건을 만족하는 행만 반환한다|	기본적으로 형태를 유지하고 값을 대체한다
|query()는 SQL이다	|Pandas의 표현식 평가 및 행 선택 메서드다|
|같은 길이의 마스크라면 언제나 안전하다|	Index 레이블이 다르면 선택 결과가 달라지거나 오류가 발생할 수 있다|
|조건부 선택은 원본을 변경한다|	조회와 할당은 별개의 연산이다|
|결측을 False로 바꾸는 것은 단순한 형식 변환이다|	판정 불가 상태를 불만족으로 처리하는 정책이다|
\end{dhktable}

여기서 \zsh{D.5}와 연결되는 핵심은 마스크의 길이만 검사해서는 충분하지 않다는 사실이다.

조건부 선택에서 중요한 것은 마스크가 어떤 행을 식별하는지, 그 Index가 원본 데이터와 올바르게 대응하는지, 결측값이 어떤 정책으로 처리되는지다.

\subsubsection{핵심 정리와 다음 학습 위치}

이번 학습의 내용을 Python 객체 모델 관점에서 정리해 보자.

\begin{dhkcode}
mask = df["score"] >= 80
result = df.loc[mask]
\end{dhkcode}

첫 번째 표현식은 Series 객체에 비교 연산을 적용해 새로운 불리언 Series를 생성한다.

두 번째 표현식은 DataFrame의 \zsh{.loc} 인덱서 객체에 불리언 Series를 전달해 행을 선택한다.

따라서 다음 흐름을 기억하면 된다.

\begin{dhksmry}[centered]
데이터 객체 \ra 조건 평가 \ra 불리언 객체 \ra 선택 또는 변환
\end{dhksmry}

그리고 Nullable Boolean이 개입한다면 이 과정에 판정 불가 상태에 대한 처리 규칙까지 추가된다.

\begin{dhknote}
\begin{dhksmry}[centered]
다음 학습 - D.7 객체 변경·재바인딩·Copy-on-Write
\end{dhksmry}
\end{dhknote}

지금까지는 주로 데이터를 조회하거나 새로운 결과 객체를 생성했다.

다음 학습에서는 기존 객체를 변경하는 과정에 집중할 것이다.

특히 다음 세 표현식의 차이를 분석할 것이다.

\begin{dhkcode}
df2 = df
df2 = df.copy()
df2 = df.assign(new_column=1)
\end{dhkcode}

그리고 다음과 같은 연쇄 대입이 왜 문제가 되는지 살펴볼 것이다.

\begin{dhkcode}
df.loc[mask, "score"] = 100
df[mask]["score"] = 100
\end{dhkcode}

이 과정에서 Python의 이름 바인딩, 객체 동일성, 내부 데이터 공유, Copy-on-Write가 서로 어떻게 연결되는지 분석하게 된다.

공식 참고 자료

\begin{dhksmry}
https://pandas.pydata.org/docs/user_guide/indexing.html
https://pandas.pydata.org/docs/user_guide/boolean.html
https://pandas.pydata.org/docs/reference/api/pandas.DataFrame.query.html
https://pandas.pydata.org/docs/reference/api/pandas.DataFrame.where.html
\end{dhksmry}